\documentclass[12pt,a4paper]{article}
\usepackage[utf8]{inputenc}
\usepackage[russian]{babel}
\usepackage{graphicx}
\usepackage{xcolor}
\usepackage{caption}
\usepackage{listings}
\usepackage{geometry}

\geometry{left=2.5cm,right=2.5cm,top=2.5cm,bottom=2.5cm}

\definecolor{MyPink}{RGB}{194, 19, 182}
\definecolor{CodeBg}{RGB}{245, 245, 245}
\definecolor{MyBlue}{RGB}{0, 102, 204}

\lstset{
    basicstyle=\ttfamily\small,
    backgroundcolor=\color{CodeBg},
    frame=single,
    breaklines=true,
    postbreak=\mbox{\textcolor{red}{$\hookrightarrow$}\space},
}

\title{\textbf{Отчет по изучению создания научных постеров\\с использованием пакета tikzposter}}
\author{Студент группы НПМмд02-24 Викторов Егор}
\date{\today}

\begin{document}

\maketitle
\tableofcontents
\newpage

\section{Введение}

Научные постеры являются важным инструментом для представления результатов исследований на конференциях, симпозиумах и других научных мероприятиях. Система LaTeX предоставляет мощный инструмент для создания профессионально оформленных постеров — пакет \texttt{tikzposter}.

\subsection{Цель работы}

Целью данной работы является изучение основных возможностей пакета \texttt{tikzposter} и освоение методов создания структурированных научных постеров.

\subsection{Задачи}

В рамках работы были поставлены следующие задачи:
\begin{itemize}
    \item Изучить работу с многоколоночными макетами
    \item Освоить создание и форматирование блоков контента
    \item Научиться вставлять графические элементы
    \item Изучить работу с цветовым оформлением
    \item Освоить добавление заметок и аннотаций
\end{itemize}

\section{Структура постера и работа с колонками}

\subsection{Основы работы с колонками}

Пакет \texttt{tikzposter} предоставляет удобное окружение \texttt{columns} для создания многоколоночных макетов. Это позволяет эффективно организовать пространство постера и разместить информацию логичным образом.

\subsubsection{Создание колонок равной ширины}

Для создания трех колонок одинаковой ширины используется следующая структура:

\begin{lstlisting}[language=TeX]
\begin{columns}
  \column{.33}
  % Содержимое первой колонки
  
  \column{.33}
  % Содержимое второй колонки
  
  \column{.33}
  % Содержимое третьей колонки
\end{columns}
\end{lstlisting}

Числовое значение в команде \texttt{\textbackslash column} определяет относительную ширину колонки. Значение 0.33 означает, что колонка занимает 33\% от общей ширины постера.

\subsubsection{Создание колонок различной ширины}

Часто требуется создать асимметричный макет с колонками разной ширины. Например, для размещения основного контента в широкой колонке и дополнительной информации в узкой:

\begin{lstlisting}[language=TeX]
\begin{columns}
  \column{.7}
  % Широкая колонка (70% ширины)
  \block{Основные результаты}{
    Здесь размещается основной контент
  }
  
  \column{.3}
  % Узкая колонка (30% ширины)
  \block{Дополнительно}{
    Здесь размещается дополнительная информация
  }
\end{columns}
\end{lstlisting}

\subsection{Практические рекомендации по работе с колонками}

\begin{enumerate}
    \item Сумма ширин колонок должна быть близка к 1.0 для оптимального использования пространства
    \item Для стандартного постера рекомендуется использовать 2-3 колонки
    \item Широкие колонки подходят для основного текста и графиков
    \item Узкие колонки удобны для списков, определений и ссылок
\end{enumerate}

\section{Работа с блоками контента}

\subsection{Базовый синтаксис блоков}

Весь контент в постере должен быть размещен внутри блоков. Блок создается с помощью команды \texttt{\textbackslash block}:

\begin{lstlisting}[language=TeX]
\block{Заголовок блока}{Содержимое блока}
\end{lstlisting}

Первый аргумент — это заголовок блока, второй — его содержимое.

\subsection{Блоки без заголовка}

Если заголовок не требуется, первые фигурные скобки оставляются пустыми:

\begin{lstlisting}[language=TeX]
\block{}{
  Это блок без заголовка. Он может содержать любой контент.
}
\end{lstlisting}

\subsection{Добавление заметок к блокам}

Одной из интересных возможностей пакета \texttt{tikzposter} является добавление заметок — небольших текстовых блоков со стрелками, указывающими на определенные части основного блока.

\subsubsection{Синтаксис команды note}

\begin{lstlisting}[language=TeX]
\block{Заголовок}{
  Основное содержимое блока
  \vspace{2cm}
}
\note[targetoffsetx=5cm, targetoffsety=-3cm, width=8cm]{
  Это заметка, которая указывает на блок выше
}
\end{lstlisting}

\subsubsection{Параметры заметок}

\begin{itemize}
    \item \texttt{targetoffsetx} — горизонтальное смещение стрелки относительно центра блока
    \item \texttt{targetoffsety} — вертикальное смещение стрелки
    \item \texttt{width} — ширина заметки
    \item \texttt{\textbackslash vspace} — создает вертикальное пространство для размещения заметки
\end{itemize}

\section{Вставка графических элементов}

\subsection{Особенности работы с изображениями}

В пакете \texttt{tikzposter} нельзя использовать стандартное окружение \texttt{figure}. Вместо этого применяется окружение \texttt{center} в сочетании с командой \texttt{\textbackslash captionof}.

\subsection{Вставка растровых изображений}

Для вставки изображений (PNG, JPG и др.) используется следующий синтаксис:

\begin{lstlisting}[language=TeX]
\block{Результаты эксперимента}{
  \begin{center}
    \includegraphics[width=0.8\textwidth]{image.png}
    \captionof{figure}{Подпись к изображению}
  \end{center}
}
\end{lstlisting}

\subsection{Вставка векторных рисунков TikZ}

Аналогичным образом вставляются рисунки, созданные с помощью пакета TikZ:

\begin{lstlisting}[language=TeX]
\block{Схема эксперимента}{
  \begin{center}
    \begin{tikzpicture}
      \draw[thick,->] (0,0) -- (4,0) node[right]{$x$};
      \draw[thick,->] (0,0) -- (0,3) node[above]{$y$};
      \draw[blue,thick] (0,0) parabola (3,2);
    \end{tikzpicture}
    \captionof{figure}{График функции}
  \end{center}
}
\end{lstlisting}

\subsection{Требования для работы с подписями}

\textbf{Важно:} Для корректной работы команды \texttt{\textbackslash captionof} необходимо подключить пакет \texttt{caption} в преамбуле документа:

\begin{lstlisting}[language=TeX]
\usepackage{caption}
\end{lstlisting}

\section{Работа с цветовым оформлением}

\subsection{Ограничения пакета tikzposter}

Пакет \texttt{tikzposter} не поддерживает опцию \texttt{svgnames} пакета \texttt{xcolor}, поэтому именованные цвета (такие как \texttt{Red}, \texttt{Blue}, \texttt{ForestGreen}) недоступны напрямую.

\subsection{Определение пользовательских цветов}

Все необходимые цвета должны быть определены в преамбуле документа с использованием команды \texttt{\textbackslash definecolor}:

\begin{lstlisting}[language=TeX]
\definecolor{MyPink}{RGB}{194, 19, 182}
\definecolor{MyBlue}{RGB}{0, 102, 204}
\definecolor{MyGreen}{RGB}{34, 139, 34}
\end{lstlisting}

\subsection{Использование цветов в тексте}

После определения цвета его можно использовать в тексте с помощью команды \texttt{\textbackslash color}:

\begin{lstlisting}[language=TeX]
Обычный черный текст, \color{MyPink} розовый текст, 
\color{black} снова черный текст.
\end{lstlisting}

\subsection{Демонстрация цветов}

Вот пример использования определенных цветов:

Это обычный черный текст. \color{MyPink}А это текст розового цвета.\color{black} Теперь снова черный. \color{MyBlue}Это синий текст.\color{black} И опять черный.

\subsection{Инструменты для подбора цветов}

Для подбора RGB-значений цветов рекомендуется использовать:
\begin{itemize}
    \item Google Color Picker — онлайн-инструмент для выбора цветов
    \item Adobe Color — для создания цветовых палитр
    \item Coolors.co — генератор цветовых схем
\end{itemize}

\section{Полный пример постера}

Ниже представлен пример структуры простого постера:

\begin{lstlisting}[language=TeX]
\documentclass[25pt, a0paper, portrait]{tikzposter}
\usepackage[utf8]{inputenc}
\usepackage[russian]{babel}
\usepackage{caption}

\definecolor{MyPink}{RGB}{194, 19, 182}

\title{Название исследования}
\author{Автор}
\institute{Университет}

\begin{document}
\maketitle

\begin{columns}
  \column{.5}
  \block{Введение}{
    Текст введения
  }
  
  \block{Методы}{
    Описание методов
  }
  
  \column{.5}
  \block{Результаты}{
    \begin{center}
      \includegraphics[width=0.9\textwidth]{results.png}
      \captionof{figure}{Основные результаты}
    \end{center}
  }
  
  \block{Выводы}{
    Основные выводы исследования
  }
\end{columns}

\end{document}
\end{lstlisting}

\section{Практические рекомендации}

\subsection{Планирование структуры}

\begin{enumerate}
    \item Определите основные разделы постера перед началом работы
    \item Спланируйте количество и ширину колонок
    \item Оцените необходимое пространство для графиков и изображений
    \item Продумайте цветовую схему
\end{enumerate}

\subsection{Оформление контента}

\begin{itemize}
    \item Используйте краткие и информативные заголовки блоков
    \item Избегайте больших объемов текста — постер должен быть визуально привлекательным
    \item Используйте списки для структурирования информации
    \item Добавляйте графики и диаграммы для наглядности
\end{itemize}

\subsection{Работа с цветами}

\begin{itemize}
    \item Используйте не более 3-4 основных цветов
    \item Обеспечьте достаточный контраст между текстом и фоном
    \item Используйте цвета для выделения ключевой информации
    \item Сохраняйте единообразие цветовой схемы по всему постеру
\end{itemize}

\subsection{Технические аспекты}

\begin{itemize}
    \item Всегда подключайте пакет \texttt{caption} для работы с подписями
    \item Используйте \texttt{\textbackslash vspace} для создания пространства под заметки
    \item Проверяйте сумму ширин колонок (должна быть ≈ 1.0)
    \item Сохраняйте изображения в высоком разрешении для печати
\end{itemize}

\section{Заключение}

В ходе выполнения данной работы были изучены основные возможности пакета \texttt{tikzposter} для создания научных постеров в системе LaTeX. 

\subsection{Освоенные навыки}

\begin{itemize}
    \item Создание многоколоночных макетов с различными пропорциями
    \item Работа с блоками контента и их заголовками
    \item Добавление заметок и аннотаций к блокам
    \item Вставка растровых и векторных изображений
    \item Определение и использование пользовательских цветов
    \item Форматирование текста и структурирование информации
\end{itemize}

\subsection{Преимущества использования tikzposter}

\begin{enumerate}
    \item Профессиональное оформление без необходимости использования графических редакторов
    \item Гибкая система компоновки контента
    \item Интеграция с другими пакетами LaTeX
    \item Возможность создания постеров различных размеров
    \item Легкость внесения изменений и обновления контента
\end{enumerate}

\subsection{Перспективы применения}

Полученные знания позволяют создавать профессионально оформленные научные постеры для:
\begin{itemize}
    \item Научных конференций
    \item Студенческих симпозиумов
    \item Защиты курсовых и дипломных работ
    \item Презентации результатов исследований
    \item Образовательных мероприятий
\end{itemize}

Пакет \texttt{tikzposter} является мощным инструментом для создания визуально привлекательных и информативных постеров, который может значительно повысить качество представления научных результатов.

\end{document}
